% Bagian: propositional-logic
% Bab: syntax-and-semantics
% Subbagian: formulas

\documentclass[../../../include/open-logic-section]{subfiles}

\begin{document}

\olfileid[id]{pl}{syn}{fml}

\olsection{\usetoken{P}{formula} Proposisi}

!!^{formula}s dalam logika proposisi dibangun dari
\emph{!!{propositional
    variable}s}\iftag{prvFalse}{\iftag{prvTrue}{,}{ dan}
  konstanta proposisional~$\lfalse$}{}\iftag{prvTrue}{\iftag{prvFalse}{
    dan}{} konstanta proposisional~$\ltrue$}{} dengan
menggunakan \emph{penghubung logis}.

\begin{enumerate}
\item Himpunan !!a{denumerable}~$\PVar$ dari !!{propositional variable}s
  $\Obj p_0$, $\Obj p_1$, \dots
\tagitem{prvFalse}{Konstanta proposisional untuk !!{falsity}~$\lfalse$.}{}
\tagitem{prvTrue}{Konstanta proposisional untuk !!{truth}~$\ltrue$.}{}
\item Penghubung logis:
  \startycommalist
  \iftag{prvNot}{\ycomma $\lnot$ (negasi)}{}%
  \iftag{prvAnd}{\ycomma $\land$ (konjungsi)}{}%
  \iftag{prvOr}{\ycomma $\lor$ (disjungsi)}{}%
  \iftag{prvIf}{\ycomma $\lif$ (!!{conditional})}{}%
  \iftag{prvIff}{\ycomma $\liff$ (!!{biconditional})}{}%
\item Tanda baca: (, ), dan koma.
\end{enumerate}

Kita menyatakan bahasa logika proposisi ini dengan $\Lang L_0$.

\iftag{defNot,defOr,defAnd,defIf,defIff,defTrue,defFalse,defEx,defAll}{%
Selain penghubung primitif yang diperkenalkan di atas, kita juga menggunakan
simbol-simbol \emph{terdefinisi} berikut:
\startycommalist
  \iftag{defNot}{\ycomma $\lnot$ (negasi)}{}%
  \iftag{defAnd}{\ycomma $\land$ (konjungsi)}{}%
  \iftag{defOr}{\ycomma $\lor$ (disjungsi)}{}%
  \iftag{defIf}{\ycomma $\lif$ (!!{conditional})}{}%
  \iftag{defIff}{\ycomma $\liff$ (!!{biconditional})}{}%
  \iftag{defFalse}{\ycomma $\lfalse$ (!!{falsity})}{}%
  \iftag{defTrue}{\ycomma $\ltrue$ (!!{truth})}{}.}{}

\begin{tagblock}{defNot,defOr,defAnd,defIf,defIff,defTrue,defFalse,defEx,defAll}
\begin{explain}
Suatu simbol terdefinisi secara resmi bukan bagian dari bahasa, tetapi
diperkenalkan sebagai singkatan informal: simbol itu memungkinkan kita
menyingkat formula yang akan menjadi cukup panjang jika kita hanya memakai
simbol primitif. Hal ini jelas merupakan suatu keuntungan. Namun, keuntungan
yang lebih besar ialah bahwa bukti menjadi lebih singkat. Jika suatu simbol
bersifat primitif, simbol itu harus ditangani secara tersendiri dalam bukti.
Oleh karena itu, semakin banyak simbol primitif, semakin panjang bukti kita.
\end{explain}
\end{tagblock}

% Simbol alternatif

\begin{intro}
Anda mungkin lebih mengenal istilah dan simbol yang berbeda dari yang kita
gunakan di atas. Teks-teks logika (dan para pengajar) lazim menggunakan
$\sim$, $\neg$, dan~!\ untuk ``negasi'', serta $\wedge$, $\cdot$, dan $\&$
untuk ``konjungsi''. Simbol yang lazim digunakan untuk ``implikasi'' ialah
$\rightarrow$, $\Rightarrow$, dan $\supset$.
\iftag{prvIff,defIff}{Simbol untuk ``biimplikasi'', ``implikasi dua arah'',
  atau ``ekuivalensi (material)'' ialah $\leftrightarrow$,
  $\Leftrightarrow$, dan $\equiv$.}{}
\iftag{prvFalse,defFalse}{Simbol $\lfalse$ disebut dengan berbagai nama,
  antara lain ``kepalsuan'', ``falsum'', ``absurditas'', atau
  ``bawah''.}{}
\iftag{prvTrue,defTrue}{Simbol $\ltrue$ disebut dengan berbagai nama,
  antara lain ``kebenaran'', ``verum'', atau ``atas''.}{}
\end{intro}

\begin{defn}[Formula]
\ollabel{defn:formulas}
Himpunan~$\Frm[L_0]$ dari \emph{!!{formula}s} logika proposisi didefinisikan
secara induktif sebagai berikut:
\begin{enumerate}
\tagitem{prvFalse}{$\lfalse$ merupakan !!{formula} atomik.}{}

\tagitem{prvTrue}{$\ltrue$ merupakan !!{formula} atomik.}{}

\item Setiap !!{propositional variable}~$\Obj p_i$ merupakan
  !!{formula} atomik.

\tagitem{prvNot}{Jika $!A$ merupakan !!a{formula}, maka $\lnot !A$
  merupakan !!a{formula}.}{}

\tagitem{prvAnd}{Jika $!A$ dan $!B$ merupakan !!{formula}s, maka
  $(!A \land !B)$ merupakan !!a{formula}.}{}

\tagitem{prvOr}{Jika $!A$ dan $!B$ merupakan !!{formula}s, maka
  $(!A \lor !B)$ merupakan !!a{formula}.}{}

\tagitem{prvIf}{Jika $!A$ dan $!B$ merupakan !!{formula}s, maka
  $(!A \lif !B)$ merupakan !!a{formula}.}{}

\tagitem{prvIff}{Jika $!A$ dan $!B$ merupakan !!{formula}s, maka
  $(!A \liff !B)$ merupakan !!a{formula}.}{}

\tagitem{limitClause}{Tidak ada ungkapan lain yang merupakan
  !!a{formula}.}{}
\end{enumerate}
\end{defn}

\begin{explain}
Definisi !!{formula}s merupakan sebuah \emph{definisi induktif}. Pada
dasarnya, kita membangun himpunan !!{formula}s dalam tak hingga banyaknya
tahap. Pada tahap awal, kita menetapkan semua formula atomik sebagai formula;
hal ini bersesuaian dengan beberapa kasus pertama dalam definisi, yaitu kasus
untuk
\iftag{prvTrue}{$\ltrue$, }{}%
\iftag{prvFalse}{$\lfalse$, }{}%
$\Obj p_i$. Dengan demikian, ``!!{formula} atomik'' berarti setiap
!!{formula} yang berbentuk demikian.

Kasus-kasus lain dalam definisi memberikan aturan untuk membangun
!!{formula}s baru dari !!{formula}s yang sudah dibangun. Pada tahap kedua,
kita dapat menggunakan aturan tersebut untuk membangun !!{formula}s dari
!!{formula}s atomik. Pada tahap ketiga, kita membangun formula baru dari
formula atomik dan formula yang diperoleh pada tahap kedua, dan demikian
seterusnya. Suatu ungkapan merupakan suatu !!{formula} jika pada akhirnya dibangun
pada salah satu tahap tersebut, dan tidak ada ungkapan lain yang merupakan
formula.
\end{explain}

Ketika menulis formula $(!B \ast !C)$ yang dibangun dari $!B$, $!C$ dengan
menggunakan penghubung dua-tempat~$\ast$, kita akan sering menghilangkan
pasangan tanda kurung terluar dan cukup menuliskan~$!B \ast !C$.

\begin{tagblock}{defNot,defOr,defAnd,defIf,defIff,defTrue,defFalse,defEx,defAll}
\begin{defn}
Formula yang dibangun dengan menggunakan operator terdefinisi harus dipahami
sebagai berikut:

\begin{tagenumerate}{defTrue,defFalse,defNot,defOr,defAnd,defIf,defIff,defEx,defAll}

\tagitem{defTrue}{$\ltrue$ menyingkat
  \iftag{prvFalse}{$\lnot\lfalse$}{$(!A \lor \lnot !A)$ untuk suatu
    !!{formula} atomik tetap~$!A$}.}{}

\tagitem{defFalse}{$\lfalse$ menyingkat
  \iftag{prvTrue}{$\lnot\ltrue$}{$(!A \land \lnot !A)$ untuk suatu
    !!{formula} atomik tetap~$!A$}.}{}

\tagitem{defNot}{$\lnot !A$ menyingkat $!A \lif \lfalse$.}{}

\tagitem{defOr}{$!A \lor !B$ menyingkat
  \iftag{prvAnd}{$\lnot(\lnot !A \land \lnot !B)$}{$\lnot !A \lif
    !B$}.}{}

\tagitem{defAnd}{$!A \land !B$ menyingkat
  \iftag{prvOr}{$\lnot(\lnot !A \lor \lnot !B)$}{$\lnot (!A \lif
    \lnot !B)$}.}{}

\tagitem{defIf}{$!A \lif !B$ menyingkat
  \iftag{prvOr}{$\lnot !A \lor !B$}{$\lnot (!A \land \lnot !B)$}.}{}

\tagitem{defIff}{$!A \liff !B$ menyingkat $(!A \lif !B) \land (!B
  \lif !A)$.}{}
\end{tagenumerate}
\end{defn}
\end{tagblock}

\begin{defn}[Identitas Sintaktis]
Simbol $\ident$ menyatakan identitas sintaktis antara untai simbol, yakni
$!A \ident !B$ jika dan hanya jika $!A$ dan $!B$ merupakan untai simbol
dengan panjang yang sama dan memuat simbol yang sama pada setiap posisi.
\end{defn}

Simbol $\ident$ dapat diapit oleh untai yang diperoleh melalui konkatenasi;
misalnya, $!A \ident (!B \lor !C)$ berarti: untai simbol~$!A$ sama dengan
untai yang diperoleh dengan menggabungkan sebuah tanda kurung buka, untai
$!B$, simbol $\lor$, untai~$!C$, dan sebuah tanda kurung tutup, dalam urutan
tersebut. Jika demikian, kita mengetahui bahwa simbol pertama dari $!A$
adalah tanda kurung buka, $!A$ memuat $!B$ sebagai subuntai (dimulai pada
simbol kedua), subuntai itu diikuti oleh $\lor$, dan seterusnya.

\end{document}
